\section{Background}
\label{sec:background}

Onion services are TCP services that are only accessible over the Tor network.
Onion domains take the place of IP addresses, which traditional TCP services use
for addressing.  These onion domains are resolved and routed inside the Tor
network, resulting in a circuit between the client and the onion service which
(by default) features six hops (see \Cref{fig:onion-connection}).

\begin{figure}[t]
\centering
\begin{tikzpicture}[node distance=0.4cm]

\tikzset{>=latex}
\tikzstyle{block} = [rectangle, draw, rounded corners, text centered,
                     minimum height=0.5cm]

\node[block,fill=black!60,
      text=white]         (TB)  {Tor Browser};
\node[block,right=of TB]  (GR1) {Guard};
\node[block,right=of GR1] (MR1) {Middle};
\node[block,below=of MR1] (R)   {Rendezvous};
\node[block,below=of R]   (MR2) {Middle};
\node[block,left=of  MR2] (MR3) {Middle};
\node[block,left=of  MR3] (GR2) {Guard};
\node[block,fill=black!60,
            text=white,
            left=of GR2]  (OS)  {Onion service};

\draw[<->] (TB.east)   -- (GR1.west);
\draw[<->] (GR1.east)  -- (MR1.west);
\draw[<->] (MR1.south) -- (R.north);
\draw[<->] (R.south)   -- (MR2.north);
\draw[<->] (MR2.west)  -- (MR3.east);
\draw[<->] (MR3.west)  -- (GR2.east);
\draw[<->] (GR2.west)  -- (OS.east);

\draw[|-|,gray] ([yshift=3pt] GR1.north west)
                -- node [above, midway, text=gray] {Circuit}
                ([yshift=3pt] MR1.north east);

\draw[|-|,gray] ([yshift=-3pt] GR2.south west)
                -- node [below, midway, text=gray] {Circuit}
                ([yshift=-3pt] MR2.south east);

\end{tikzpicture}
\caption{A connection to an onion service typically consists of six Tor relays.
Both the client and the onion service create a circuit (consisting of two and
three relays, respectively) to the rendezvous relay that serves as a short-lived
data exchange point.}
\label{fig:onion-connection}
\end{figure}

The creation of a new onion domain requires a Tor daemon to generate an RSA key
pair.  It then computes the SHA-1 hash over the RSA public key, truncates it to
80 bits, and encodes these 80 bits in Base32, resulting in sixteen characters,
\eg, expyuzz4wqqyqhjn.  As of January 2018, The Tor Project is deploying the
next generation of onion services whose domain format will feature 56 instead of
sixteen characters~\cite[\S~6]{Mathewson2013a}---a Base32 encoding of the onion
service's public key, a checksum, and a version number.  Because of the next
generation using elliptic curve cryptography, the entire public key (instead of
just a hash over the public key) is embedded in the domain.

Due to an onion domain being a function of its public key, onion domains are
self-authenticating, \ie, as long as a client has the correct domain, it knows
what public key to expect.  The downside is that sixteen random characters are
impractical to remember, let alone 56 characters.

We can make onion domains at least partially meaningful by repeatedly creating
RSA keys until the resulting domain contains a desired string.  We call these
\emph{vanity onion domains}.  A vanity prefix of length $n$ takes on average
$0.5 \cdot 32^n$ key creations given Base32's alphabet size of 32 characters.
After having created a set of domains featuring a vanity prefix, one can search
this set for the domain that is the easiest to remember, \eg, by using a Markov
model to filter domains that resemble words in the English language.  While this
method will not produce fully-meaningful domains, it can facilitate memorization
as evidenced by the vanity domains of Facebook
(\href{https://www.facebookcorewwwi.onion}{facebookcorewwwi.onion}),
ProPublica
(\href{https://www.propub3r6espa33w.onion}{propub3r6espa33w.onion}),
and the New York Times
(\href{https://www.nytimes3xbfgragh.onion}{nytimes3xbfgragh.onion}).  In
practice, many onion service operators use the tool scallion~\cite{scallion} to
find vanity domains in parallel.

Onion services are private by default.  Once an onion service is created, it is
its operator's task to disseminate the domain, \eg, by adding it to onion site
search engines such as Ahmia~\cite{ahmia}.  The lack of a go-to service such as
Google for onion service discovery prompted the community to devise various ways
to disseminate onion services, most importantly an array of search engines and
curated lists.

On the usability front, Tor Browser's user interface when accessing onion
services is designed to be seamless.  \Cref{fig:non-onion-service} shows the UI
when accessing The Tor Project's web site while \Cref{fig:onion-service} shows a
connection to the corresponding onion site.

\begin{figure}[t]
    \centering

    \begin{subfigure}[t]{\linewidth}
        \centering
        \includegraphics[width=\linewidth]{figures/non-onion-service.png}
        \subcaption{The Tor Project's web site when accessed over its
        conventional domain.}
        \label{fig:non-onion-service}
    \end{subfigure}

    \begin{subfigure}[t]{\linewidth}
        \centering
        \includegraphics[width=\linewidth]{figures/onion-service.png}
        \subcaption{The Tor Project's web site when accessed over its onion
        service.}
        \label{fig:onion-service}
    \end{subfigure}

    \caption{Tor Browser 7.0.10's user interface on Windows 10 when opening a
    conventional domain (top) and when opening the corresponding onion service
    (bottom).  Note that the onion version lacks a padlock icon---an issue
    that is currently being worked on~\cite{trac23247}.}
\end{figure}
